\bookchapter{第四章：安裝同管理App——手機嘅「技能包」}

\chapterintro{App等手機無所不能，但都要管好佢哋。就好似你會畀新員工安排工位、分配權限一樣，安裝App都需要一套規範流程。}

\booksection{一、官方商店：最安全嘅「應用超市」}

\begin{itemize}
  \item \textbf{安卓}：Google Play（國外）或各品牌自帶商店（如華為應用市場、小米應用商店）。
  \item \textbf{iOS}：淨係得 App Store。
  \item \textbf{鴻蒙}：華為應用市場。
\end{itemize}

核心安全意義：經過審核、更新有保障、支付安全。行動指南：\textbf{優先淨係用官方商店下載！}

\booksection{二、純淨/防護模式：開啟手機嘅「防毒牆」}

主流國產品牌（華為、小米、OPPO、vivo、榮耀）都內置咗\textbf{純淨模式}或\textbf{防護模式}，呢個係系統級別嘅安全開關。\textbf{強烈建議所有用戶都開啟}。

\booksubsection{佢到底係乜嘢？}

純淨模式（又叫「防護模式」）嘅\textbf{核心功能係阻斷第三方 App 嘅安裝}：當佢開啟嗰陣，系統會\textbf{淨係允許由官方應用商店安裝 App}，由瀏覽器、二維碼、聊天軟件等渠道下載嘅安裝包會被直接攔截，由源頭杜絕惡意軟件。

\booksubsection{適用人羣}

\begin{itemize}
  \item \textbf{老年用戶}：分唔清「官方」同「假冒」廣告，純淨模式係最穩妥嘅護身符。
  \item \textbf{對應用權限唔敏感嘅用戶}：唔需要安裝小眾 App、測試版、外服遊戲。
  \item \textbf{希望提升設備安全性嘅所有用戶}：哪怕係發燒友，開住都好過唔開——真係有臨時需要，再單獨關一下即可。
\end{itemize}

\booksubsection{開啟路徑（以華為/小米為例）}

\begin{itemize}
  \item \textbf{華為/鴻蒙}：[設定] > [系統同更新] >\textbf{[純淨模式]}，選擇「開啟」。
  \item \textbf{小米}：[設定] > [私隱保護] >\textbf{[保護私隱]}或 [設定] > [應用設定] > [應用安裝]，關閉「安裝未知應用」。
  \item \textbf{OPPO/vivo}：[設定] > [安全] >\textbf{[安裝外部來源應用]}，逐個關閉授權。
\end{itemize}

\booksubsection{應用商店嘅安全檢測機制}

應用商店唔係「擺設」，佢背後有完整嘅安全流水線：

\begin{itemize}
  \item \textbf{病毒掃描}：每個 App 上架前會經過多家殺毒引擎聯合掃描，識別已知木馬同病毒特徵。
  \item \textbf{權限審核}：商店會審查 App 申請嘅權限係咪「合理」——一個手電筒 App 要通訊錄權限？直接駁回或下架。
  \item \textbf{開發者認證}：上架需要企業資質、個人身份證或信用卡認證。出問題可以揾到人，唔係「黑 App」。
  \item \textbf{持續監控}：已上架嘅 App 都會被持續複掃，新發現嘅惡意行為會被強制下架。
\end{itemize}

\booksection{三、第三方來源與「電子寄生蟲」：風險極高！}

由網頁或不明連結下載App，可能夾帶病毒、竊取私隱、盜刷支付。呢度我哋特別提一個新概念：\textbf{「電子寄生蟲」}。

\textbf{乜嘢係「電子寄生蟲」？}\par

指嘅係嗰啲\textbf{被植入木馬嘅「破解版」軟件}——你以為係「慳咗錢」，實際係「中咗招」。常見形式：

\begin{itemize}
  \item 「破解版」、「免費版」、「去廣告版」付費軟件；
  \item 「刷讚」、「搶紅包」、「遊戲外掛」等灰色工具；
  \item 羣聊入面「小廣告」推薦嘅所謂「內部版」App。
\end{itemize}

呢啲程式表面能用，背後可能喺度\textbf{偷通訊錄、讀短訊、監聽麥克風、後台挖礦}。純淨模式係防範佢哋嘅\textbf{第一道防線}——直接唔等你裝。

\booksubsection{第二道防線建議：}

\begin{itemize}
  \item \textbf{定期全盤病毒掃描}：使用手機自帶安全中心（如華為手機管家、小米安全中心）每月掃一次。
  \item \textbf{謹慎授予敏感權限}：通訊錄、短訊、位置、相機、麥克風——呢五項「高危權限」必須嚴格審批。
  \item \textbf{關注耗電同流量異常}：手機突然發燙、電量跌得快、流量半夜跑——八成係後台收埋咗「寄生蟲」。
\end{itemize}

如果萬不得已要裝第三方 App：

\begin{enumerate}
  \item 確保來源絕對可信。
  \item 安裝前仔細睇權限要求。
  \item 事後用安全App掃描。
\end{enumerate}

終極建議：喺系統設定入面，關掉「安裝未知應用」嘅權限。

\booksection{四、裝、更新、卸載}

\begin{itemize}
  \item \textbf{安裝}：商店入面點「取得」或「安裝」，裝完首次打開會問你要權限。
  \item \textbf{更新}：\textbf{必須更！}主要係修復安全漏洞。建議開啟自動更新。
  \item \textbf{卸載}：長按圖標拖到卸載區。諗住清得更徹底，去 [設定] > [應用管理] 入面揾到並卸載。
\end{itemize}

\booksection{五、權限管理：你係App嘅「老闆」}

App要權限（比如通訊錄、位置）有時合理（地圖要位置），有時唔合理（手電筒要通訊錄）。

\begin{itemize}
  \item \textbf{權限選項}：「始終允許」、「僅喺使用時允許」、「禁止」、「每次詢問」。
  \item \textbf{管理原則}：
  \begin{enumerate}
    \item 最小化：非必要唔畀。
    \item 事後審計：定期去權限設定入面睇下邊個喺度「越權」。
    \item 情境化：好似打車軟件，淨係用嗰陣畀位置權限。
  \end{enumerate}
\end{itemize}

\begin{tcbwarn}{重點盯防嘅權限}
\begin{itemize}
  \item 通訊錄/短訊：可能被用嚟發垃圾短訊。
  \item 實時位置：可能被追蹤行蹤。
  \item 相機/麥克風：可能被偷拍偷錄。
  \item 檔案訪問：可能讀取你嘅私隱檔案。
\end{itemize}

\end{tcbwarn}

\booksection{六、軟件安裝包嘅構成要素}

你有冇諗過，由應用商店下載嗰個安裝包，入面到底裝咗乜嘢？瞭解安裝包嘅構成，可以幫你更好咁理解App係點樣工作嘅，都可以喺遇到問題嗰陣知道該由邊度排查。

\booksubsection{安卓APK檔案結構}

APK（Android Package Kit）本質係一個ZIP壓縮包，解壓後可以睇到以下核心目錄同檔案：

\begin{itemize}
  \item \textbf{AndroidManifest.xml}：應用嘅「身份證」同「說明書」。記錄咗應用包名、版本號、所需權限、支援嘅系統版本、入口Activity等核心資訊。系統安裝應用嗰陣首先讀取呢個檔案。
  \item \textbf{classes.dex}：應用嘅核心代碼。Java/Kotlin源代碼編譯後變成Dalvik字節碼，存喺呢個檔案入面。一個應用可能有多個dex檔案（如classes2.dex、classes3.dex）。
  \item \textbf{resources.arsc}：編譯後嘅資源索引表。記錄咗所有資源嘅ID同對應關係。
  \item \textbf{res/目錄}：資源資料夾。包含圖片、佈局檔案、字串、顏色定義等。唔同分辨率嘅圖片放喺唔同嘅drawable資料夾入面（如drawable-hdpi、drawable-xhdpi）。
  \item \textbf{assets/目錄}：原始資源資料夾。存放唔經過編譯嘅資源，如網頁檔案、字體、音頻等。App運行嗰陣可以通過AssetManager讀取。
  \item \textbf{lib/目錄}：本地庫資料夾。存放編譯好嘅.so檔案（C/C++代碼編譯產物），按CPU架構分類（如arm64-v8a、armeabi-v7a）。遊戲同需要高性能計算嘅App通常有呢個目錄。
  \item \textbf{META-INF/目錄}：簽名資訊資料夾。包含MANIFEST.MF、CERT.SF、CERT.RSA三個檔案，用於驗證APK嘅完整性同簽名。
\end{itemize}

\booksubsection{iOS IPA檔案結構}

IPA（iOS App Store Package）同樣係壓縮包，結構略有唔同：

\begin{itemize}
  \item \textbf{Payload/目錄}：核心內容所在。入面有一個.app資料夾，包含應用嘅所有內容。
  \item \textbf{Info.plist}：應用嘅配置檔案，類似安卓嘅AndroidManifest.xml。記錄咗應用ID、版本、權限需求、支援嘅設備等。
  \item \textbf{可執行檔案}：與.app資料夾同名嘅主程式，係編譯後嘅Mach-O格式二進制檔案。
  \item \textbf{Frameworks/目錄}：應用依賴嘅動態庫同框架。
  \item \textbf{資源檔案}：圖片、音頻、影片等資源，通常存放在.app資料夾內。
  \item \textbf{\_CodeSignature/目錄}：代碼簽名資訊，用於驗證應用嘅完整性同來源。
\end{itemize}

\booksubsection{簽名驗證機制}

點解由官方商店下載嘅App更安全？簽名驗證係關鍵。

\begin{itemize}
  \item \textbf{安卓簽名}：開發者用私鑰對APK進行簽名，系統安裝嗰陣用公鑰驗證。如果APK被篡改（比如植入病毒），簽名就會失效，系統會拒絕安裝。Google Play仲要求App Signing，由Google保管簽名密鑰，進一步保護應用安全。
  \item \textbf{iOS簽名}：蘋果嘅審核更嚴格。每個IPA必須有蘋果頒發嘅證書簽名，否則無法安裝。呢個就係點解iOS淨係可以由App Store或企業證書安裝應用。企業證書簽名嘅應用有效期有限，且蘋果可以隨時吊銷。
  \item \textbf{簽名嘅作用}：
  \begin{enumerate}
    \item 驗證來源：確認應用確實來自聲稱嘅開發者。
    \item 驗證完整性：確保應用喺傳輸過程中未被篡改。
    \item 防止冒充：惡意應用無法偽裝成其他應用。
  \end{enumerate}
\end{itemize}

\booksubsection{安裝包大小與優化}

你可能會發現，應用商店顯示嘅App大小同安裝後佔用嘅空間唔一樣。呢個係因為：

\begin{itemize}
  \item \textbf{下載大小}：壓縮後嘅APK/IPA大小。
  \item \textbf{安裝大小}：解壓後嘅大小，通常比下載大小大好多。
  \item \textbf{運行時佔用}：App運行嗰陣仲會產生快取、數據檔案，佔用更多空間。
\end{itemize}

現代App嘅優化手段：

\begin{itemize}
  \item \textbf{App Bundle（安卓）}：Google Play嘅動態分發技術。用戶下載嗰陣淨係獲取適合自己設備嘅資源（如淨係下載對應螢幕分辨率嘅圖片），大大減少下載大小。
  \item \textbf{App Thinning（iOS）}：類似嘅技術，蘋果會根據設備自動裁剪唔需要嘅資源。
  \item \textbf{按需下載}：遊戲等大型應用可以先下載核心包，後續再下載資源包。
\end{itemize}

\begin{tcbinfo}{小知識：點解有啲App安裝包特別大？}
遊戲類App通常最大，因為包含大量高清貼圖、3D模型、音頻檔案。社交類App都比較大，因為要支援各種功能模組。工具類App通常最小。如果你發現某個工具類App特別大，要警惕係咪捆綁咗唔必要嘅功能。

\end{tcbinfo}

\begin{tcbpractice}{實操環節：審查3個App嘅權限}
揀3個常用App（如社交、購物、工具類），進入 [設定] 嘅應用權限管理頁，逐項審查，將唔必要嘅敏感權限關掉或改成「僅喺使用時允許」。

\end{tcbpractice}
